\documentclass[10pt,a4paper]{article}
\usepackage[a4paper,top=17mm,bottom=20mm,left=16.5mm,right=16.5mm,headheight=13pt,headsep=5mm,footskip=10mm]{geometry}
\usepackage{fontspec}
\usepackage{amsmath,amssymb,mathtools,booktabs,tabularx,array,multicol,caption}
\usepackage{xcolor,hyperref,fancyhdr,lastpage,orcidlink,tikz}
\usetikzlibrary{positioning}
\definecolor{ink}{HTML}{111111}
\definecolor{muted}{HTML}{5B6168}
\definecolor{lightbar}{HTML}{888888}
\definecolor{rulec}{HTML}{BBBBBB}
\setmainfont{Norasi.otf}[Ligatures=TeX,BoldFont=Norasi-Bold.otf,ItalicFont=Norasi-Italic.otf,BoldItalicFont=Norasi-BoldItalic.otf]
\newfontfamily\uisans{Garuda.otf}[Ligatures=TeX,BoldFont=Garuda-Bold.otf,ItalicFont=Garuda-Oblique.otf,BoldItalicFont=Garuda-BoldOblique.otf]
\XeTeXlinebreaklocale "th"
\XeTeXlinebreakskip=0pt plus .65pt
\emergencystretch=0.7em
\setlength{\parindent}{1em}
\setlength{\parskip}{0pt}
\setlength{\columnsep}{6.4mm}
\setlength{\abovedisplayskip}{6pt plus2pt minus2pt}
\setlength{\belowdisplayskip}{6pt plus2pt minus2pt}
\captionsetup{font=small,labelfont=bf,labelsep=period}
\hypersetup{
  colorlinks=true, linkcolor=ink, citecolor=ink, urlcolor=ink,
  pdftitle={การจัดพิมพ์วิชาการภาษาไทยระดับสากลด้วย LaTeX - Scientific Letter},
  pdfauthor={Yaoharee Lahtee},
  pdfsubject={Two-column Thai scientific typesetting},
  pdfkeywords={Thai; XeLaTeX; two-column; scientific letter; typography}
}
\newcommand{\TB}{\penalty0\hskip0.14em plus0.025em minus0.012em\relax}
\newcommand{\MetaSep}{\hspace{0.45em}\textbullet\hspace{0.45em}}
\newcommand{\shorttitle}{THAI SCIENTIFIC TYPESETTING WITH LATEX}
\fancypagestyle{journal}{%
  \fancyhf{}
  \fancyhead[L]{\uisans\fontsize{7.3}{8.5}\selectfont YAOHAREE LAHTEE}
  \fancyhead[C]{\uisans\fontsize{7.3}{8.5}\selectfont\shorttitle}
  \fancyhead[R]{\uisans\fontsize{7.3}{8.5}\selectfont PREPRINT}
  \renewcommand{\headrulewidth}{0.45pt}
  \fancyfoot[L]{\uisans\fontsize{7.1}{8.3}\selectfont OCSI-TH-2026-001}
  \fancyfoot[C]{\uisans\fontsize{7.1}{8.3}\selectfont \thepage\ / \pageref*{LastPage}}
  \fancyfoot[R]{\uisans\fontsize{7.1}{8.3}\selectfont CC BY 4.0 \MetaSep 10.XXXX/ocsi.2026.001}
  \renewcommand{\footrulewidth}{0.25pt}
}
\pagestyle{journal}
\begin{document}
\noindent\rule{\textwidth}{0.65pt}\par\vspace{4pt}
\noindent\colorbox{lightbar}{\makebox[31mm][l]{\color{white}\uisans\bfseries\fontsize{7.8}{9}\selectfont LETTER / METHODS}}\hfill
{\uisans\fontsize{7.5}{9}\selectfont OPEN ACCESS \MetaSep PREPRINT}\par
\vspace{8pt}
\begin{center}
{\uisans\bfseries\fontsize{17.6}{20.4}\selectfont การจัดพิมพ์วิชาการภาษาไทยระดับสากลด้วย LaTeX:\par
ระบบสองคอลัมน์สำหรับสมการ ตาราง และ metadata\par}
\vspace{4pt}
{\fontsize{9.8}{11.6}\selectfont \textbf{Yaoharee Lahtee}\,\orcidlink{0009-0005-3861-0626}\,*\par}
{\fontsize{8.7}{10.4}\selectfont Open Civil Science Initiative, Bangkok, Thailand\par}
{\uisans\fontsize{7.6}{9}\selectfont (Received 19 September 2026; revised --; accepted --; version 1.0)}
\end{center}
\vspace{4pt}
\begin{center}\begin{minipage}{0.94\textwidth}
\fontsize{9.15}{11.7}\selectfont\setlength{\emergencystretch}{1.5em}
ระบบจัดพิมพ์เชิงวิทยาศาสตร์แบบสองคอลัมน์เพิ่มความหนาแน่นของข้อมูล\TB แต่สร้างข้อจำกัดต่อสมการยาว\TB ตาราง\TB รูป\TB และข้อความสองภาษาโดยตรง\TB งานนี้เสนอ workflow ที่ให้โครงสร้างมาก่อน scaling\TB อนุญาตให้วัตถุขยายเต็มสองคอลัมน์เมื่อจำเป็น\TB และใช้ preflight แบบ fail-closed เพื่อตรวจ overflow\TB font embedding\TB Unicode text layer\TB และ metadata ก่อนเผยแพร่
\par\smallskip
\noindent\textbf{DOI:} 10.XXXX/ocsi.2026.001 \qquad
\textbf{arXiv:} XXXX.XXXXX \qquad
\textbf{ORCID:} 0009-0005-3861-0626
\end{minipage}\end{center}
\vspace{4pt}\noindent\rule{\textwidth}{0.55pt}\vspace{5pt}

\begin{multicols}{2}
\raggedcolumns
\fontsize{9.15}{11.7}\selectfont
\noindent\textbf{บทนำ.}\quad
เอกสารสองคอลัมน์ควรทำให้ผู้อ่านสแกนโครงสร้างของงานได้รวดเร็ว\TB โดยไม่ทำให้คอลัมน์แคบจนสมการและตารางเสียรูป\TB การตัดสินใจเรื่อง layout จึงต้องทำที่ระดับโครงสร้าง\TB ไม่ใช่ลดขนาดตัวอักษรจนอ่านยาก

\medskip
\noindent\textbf{Semantic composition.}\quad
ภาษาไทยควรถูกแบ่งตามหน่วยคิดและน้ำหนักการหายใจก่อนเข้าสู่ compositor\TB ความยาวของหน่วยไม่ควรเท่ากันเป็นจังหวะเครื่องจักร\TB และการจัดเต็มบรรทัดต้องไม่แลกกับช่องว่างที่ผิดธรรมชาติ

\begin{equation}
Q_{\rm pub}=Q_{\rm sem}\,Q_{\rm layout}\,Q_{\rm PDF}.
\end{equation}

\noindent\textbf{Column-aware math.}\quad
สมการที่อยู่ได้ในหนึ่งคอลัมน์ควรรักษาขนาดตัวอักษรปกติ\TB หากยาวเกินควรใช้ aligned หรือ split ก่อน\TB แล้วจึงพิจารณาขยายเต็มสองคอลัมน์

\medskip
\noindent\textbf{Metadata.}\quad
ORCID เชื่อมตัวตนผู้เขียน\TB DOI และ arXiv identifier เชื่อมระเบียนบทความ\TB ส่วน version และ article history ช่วยให้การอ้างอิงฉบับที่ถูกต้องทำได้โดยไม่ต้องเดา

\medskip
\noindent\textbf{PDF integrity.}\quad
การ compile สำเร็จยังไม่เพียงพอ\TB pipeline ต้องตรวจ missing glyph\TB overfull boxes\TB ฟอนต์ฝัง\TB ToUnicode\TB unresolved references\TB และ text extraction ของภาษาไทยด้วย

\columnbreak
\noindent\textbf{Layout escalation.}\quad
ลำดับการแก้ควรเริ่มจาก recompose และ reflow\TB ต่อด้วย full-width object\TB landscape หรือ appendix ตามชนิดของเนื้อหา\TB scaling เป็นเพียงทางเลือกสุดท้ายและต้องมี threshold ที่ระบบสามารถปฏิเสธได้

\medskip
\noindent\textbf{Typography.}\quad
สองคอลัมน์ระดับวารสารต้องรักษา baseline\TB page color\TB และความยาวบรรทัดให้สม่ำเสมอ\TB แต่ไม่ควรบังคับให้ทุกคอลัมน์มีจำนวนบรรทัดเท่ากัน หากทำให้ float หรือสมการผิดลำดับเหตุผล

\medskip
\noindent\textbf{Equation policy.}\quad
สมการควรถูกแตกตามตรรกะของนิพจน์ก่อน\TB ตัวห้อยและเศษส่วนต้องไม่เล็กจนต่ำกว่าระดับอ่านสบาย\TB และ full-width equation ต้องยังรักษาหมายเลขไว้ในแนวเดียวกับสมการปกติ

\medskip
\noindent\textbf{Table policy.}\quad
ตารางควรใช้ booktabs และพื้นที่ reflow ก่อน\TB หลีกเลี่ยงเส้นกรอบทุกช่อง\TB และย้ายไป full width หรือ appendix เมื่อข้อมูลไม่เหมาะกับ column width

\medskip
\noindent\textbf{Correspondence.}\quad
* Corresponding author: [research email].\TB ORCID: 0009-0005-3861-0626.\TB License: CC BY 4.0.\TB Preprint / not peer reviewed.
\end{multicols}

\vspace{5pt}
\begin{center}
{\uisans\bfseries\fontsize{8.2}{9.8}\selectfont FULL-WIDTH MATHEMATICAL OBJECT}\par\vspace{3pt}
\begin{equation}
\mathcal{K}_A(W)=1
\quad\Longleftrightarrow\quad
\sup_{\nu\in\mathcal{V}(W)}
\operatorname{Dist}\!\left(R_A^{\nu}[n],R_A[n]\right)
\le \varepsilon_K
\quad\text{for all admissible }\nu\text{ on }W.
\label{eq:wide}
\end{equation}
\end{center}
\vspace{4pt}

\begin{center}
\captionof{table}{Policy สำหรับวัตถุที่มีความเสี่ยงต่อ overflow ในเอกสารสองคอลัมน์}\label{tab:widepolicy}
\small
\begin{tabularx}{0.94\textwidth}{@{}>{\bfseries}lXXl@{}}
\toprule
ประเภท & ขั้นที่ 1 & ขั้นที่ 2 & Fail-closed trigger\\
\midrule
สมการยาว & split / aligned / multline & full-width equation & width ก่อนย่อเกิน 115\%\\
ตารางกว้าง & reflow / tabularx & full width / landscape / appendix & scale รุนแรงหรืออ่านไม่ได้\\
รูปกว้าง & crop อย่างมีเหตุผล & span สองคอลัมน์ & label หรือ caption ถูกตัด\\
ภาษาไทย & semantic recompose & controlled break opportunity & spacing ผิดธรรมชาติ\\
\bottomrule
\end{tabularx}
\end{center}
\end{document}
